\documentclass[10pt]{article}

% Specify the margins and text size.
\setlength{\textwidth}{6.4in}
\setlength{\textheight}{9.5in}
\setlength{\oddsidemargin}{0pt}
\setlength{\evensidemargin}{0pt}
\setlength{\topmargin}{0pt}
\setlength{\hoffset}{.05in}
\setlength{\voffset}{-.8in}

\setlength{\parskip}{5pt}
\setlength{\parindent}{0pt}

% Load some fonts and symbol packages
\usepackage{latexsym}
\usepackage{pifont}       % contains 'star' symbol for counterinsurgency handbook title
\usepackage{yfonts} 
\usepackage{amsmath}
\usepackage{amsfonts}

\usepackage{graphicx}     % actually, this is loaded by pstricks
\usepackage[T1]{fontenc}
\usepackage{ifthen}
\usepackage{pstricks,pst-grad,pst-text,pst-node,multido,pst-plot,calc,pst-3dplot}
%\usepackage[all]{xy}
%\usepackage{animate}

% The hyperref package inserts links into the pdf.
\definecolor{MyLinkColor}{rgb}{.1,.2,1}
\definecolor{MyCiteColor}{rgb}{.1,1,.2}
\definecolor{MyURLColor}{rgb}{.2,.2,.6}
\usepackage[backref=true,pagebackref=false,hyperindex,colorlinks=true,
  linkcolor=MyLinkColor,urlcolor=MyURLColor]{hyperref}


% The tweaklist package is something I found on the web.  It provides a simple interface
% for making changes to spacing used in the itemize and enumerate environments.  Comment
% this out if you don't care to use tweaklist.
\usepackage{tweaklist}
\renewcommand{\itemhook}{\setlength{\parskip}{2pt}\setlength{\parsep}%
{1pt}\setlength{\topsep}{0pt}\setlength{\itemsep}{0pt}}

\newcommand{\U}{\underline{\hspace{5pt}}}

\usepackage{listings}
\newcounter{EX}\setcounter{EX}{1}
\newcommand{\EXERCISE}{\arabic{EX}.\stepcounter{EX} }

\begin{document}
\pagestyle{empty}
\lstset{language=R, showspaces=false, showstringspaces=false}

\href{http://www.su.edu}{\includegraphics[height=1.75cm]{sulogo.eps}}
\vspace{-1.69cm}

{\small \hfill
\begin{tabular}{cl}
& Math 207\\
& Introduction to Statistics\\
\end{tabular}
}
\setlength{\baselineskip}{1.05\baselineskip}

\begin{center}
\textbf{\large  Installing \texttt{R} and RStudio}
\end{center}
\medskip

The \texttt{R} Environment for Statistical Computing 
 is a free software environment for statistical computing and graphics.
RStudio is a graphical user interface to \texttt{R}.  
Both are used extensively in academics, industry and government.
They are freely-available to download for Mac, Linux and Windows.
We will use both extensively this semester.  Please follow the steps below
to install and test them on your computer.
\medskip


\textbf{Install R}:

\begin{itemize}
\item Open a web browser to \href{https://www.r-project.org}{https://www.r-project.org}.
\item Click on the \textit{download R} link near the top of the page.
\item Scroll down the CRAN Mirrors page and select a U.S. site.
\item Click on the \textit{Download for (Mac) OS X} link (or the link that is
   appropriate to your computer's operating system).
\item Once the download is complete, install the software.  On Mac, this involves
  double-clicking on the downloaded image file.
\end{itemize}
\medskip

\textbf{Test Your R Installation}:
\begin{itemize}
\item Double-click on the new R icon in your Applications menu.
\item Once R has started, type \texttt{2+3} in the window then hit return.
\item If the output is 5, then R is installed correctly.
\end{itemize}
\medskip

\textbf{Install R Studio}:
\begin{itemize}
\item Open a web browser to \href{https://www.rstudio.com}{https://www.rstudio.com}.
\item Click on the \textit{Download RStudio} link near the top of the page.
\item Click on the RStudio Desktop, Open Source License, FREE Download link.
\item Click download for your computer's operating system (e.g., Mac OSX).
\item Once the download is complete, install the software.  On Mac, this involves
  double-clicking on the downloaded image file.
\end{itemize}
\medskip


\textbf{Test Your RStudio Installation}:
\begin{itemize}
\item Double-click on the new RStudio icon in your Applications menu.
\item Once RStudio has started, type \texttt{2+3} in the command window then hit return.
\item If the output is 5, then RStudio is installed correctly.
\end{itemize}
\medskip

\textbf{Install MacTeX}:
\begin{itemize}
\item Open a web browser to \href{http://www.tug.org/mactex}{http://www.tug.org/mactex}.
\item Click on the \textit{MacTeX Download} link near the top of the page.
\item Click on the \textit{MacTeX.pkg} link to begin the (almost 4GB!) download.
\item Once the download is complete, install the software.  On Mac, this involves
  double-clicking on the downloaded image file.
\item We won't use \TeX\ directly.  It is used by an R package to generate pdf 
   versions of reports.
\end{itemize}
\medskip

\vfill
\eject

\end{document}
